% --- UNIVERSAL PREAMBLE BLOCK ---
\documentclass[11pt, a4paper]{article}

\usepackage[utf8]{inputenc}
\usepackage[T1]{fontenc}
\usepackage[spanish, es-tabla, es-nodecimaldot]{babel}
\usepackage[a4paper, top=2.5cm, bottom=2.5cm, left=2.5cm, right=2.5cm]{geometry}

\usepackage{graphicx}
\usepackage{float}
\usepackage{hyperref}
\usepackage{enumitem}
\usepackage{xcolor}
\usepackage{listings}
\usepackage{fancyhdr}
\usepackage{tabularx}
\usepackage{longtable}
\usepackage{xltabular}
\usepackage{booktabs}
\usepackage{colortbl}
\usepackage{caption}
\usepackage{amssymb}

\setlist[itemize]{label=-}
\setcounter{secnumdepth}{0}

\definecolor{codegreen}{rgb}{0,0.6,0}
\definecolor{codegray}{rgb}{0.5,0.5,0.5}
\definecolor{codepurple}{rgb}{0.58,0,0.82}
\definecolor{backcolour}{rgb}{0.96,0.96,0.96}

\lstdefinestyle{pythonstyle}{
    backgroundcolor=\color{backcolour},
    commentstyle=\color{codegreen},
    keywordstyle=\color{magenta},
    numberstyle=\tiny\color{codegray},
    stringstyle=\color{codepurple},
    basicstyle=\ttfamily\small,
    breakatwhitespace=false,
    breaklines=true,
    captionpos=b,
    keepspaces=true,
    numbers=left,
    numbersep=5pt,
    showspaces=false,
    showstringspaces=false,
    showtabs=false,
    tabsize=4,
    language=Python,
    extendedchars=true,
    literate={á}{{\'a}}1 {é}{{\'e}}1 {í}{{\'i}}1 {ó}{{\'o}}1 {ú}{{\'u}}1 {ñ}{{\~n}}1 {Á}{{\'A}}1 {É}{{\'E}}1 {Í}{{\'I}}1 {Ó}{{\'O}}1 {Ú}{{\'U}}1 {Ñ}{{\~N}}1 {π}{{$\pi$}}1 {²}{{$^2$}}1 {³}{{$^3$}}1 {¿}{{?`}}1
}
\lstset{style=pythonstyle}

\hypersetup{
    colorlinks=true,
    linkcolor=black,
    filecolor=magenta,
    urlcolor=blue,
    pdftitle={Actividad 2 - POO},
}

% Comando auxiliar para dibujar un cuadro de clase estilo UML
\newcommand{\clasebox}[3]{%
\begin{center}
\begin{tabular}{|l|}
\hline
\textbf{#1} \\
\hline
#2
\hline
#3
\hline
\end{tabular}
\end{center}
}

\begin{document}

% ---------------------------------------------------------
% PORTADA
% ---------------------------------------------------------
\begin{titlepage}
    \centering
    \vspace*{1cm}

    {\huge \textbf{Universidad Nacional de Colombia}\par}
    \vspace{0.5cm}
    {\Large Facultad de Minas\par}
    \vspace{1.5cm}

    {\Large \textbf{Programación Orientada a Objetos 2026-2}\par}
    \vspace{2cm}

    {\huge \textbf{Actividad 2}\par}
    \vspace{0.5cm}
    {\large Seguimiento y solución de ejercicios prácticos de Programación Orientada a Objetos\par}

    \vspace{3cm}

    \textbf{Presentado por:}\par
    \vspace{0.5cm}
    {\Large \textbf{Natalia Patricia Remolina Rodríguez}\par}

    \vspace{2cm}

    \textbf{Docente:}\par
    Walter Hugo Arboleda Mazo\par

    \vspace{2cm}

    Medellín, Colombia\par
    8 de octubre de 2026
\end{titlepage}

% ---------------------------------------------------------
% EJERCICIO 2.1
% ---------------------------------------------------------
\section{Ejercicio 2.1}

\noindent \textbf{Enlace al repositorio en GitHub:} \\
\url{https://github.com/NatoRemolina/UNAL-Programacion-Orientada-a-Objetos/blob/main/Actividad%202/ejercicio1.py}

\subsection{Código fuente}
\begin{lstlisting}
class Persona:
    def __init__(self, nombre, apellido, numero_documento, año_nacimiento):
        self.nombre = nombre
        self.apellido = apellido
        self.numero_documento = numero_documento
        self.año_nacimiento = año_nacimiento

    def mostrar_datos_personales(self):
        print(f"Nombre = {self.nombre}")
        print(f"Apellido = {self.apellido}")
        print(f"Número de documento de identidad = {self.numero_documento}")
        print(f"Año de nacimiento = {self.año_nacimiento}")

nombre = input("Ingrese el nombre: ")
apellido = input("Ingrese el apellido: ")
numero_documento = input("Ingrese el número de documento de identidad: ")
año_nacimiento = int(input("Ingrese el año de nacimiento: "))

persona1 = Persona(nombre, apellido, numero_documento, año_nacimiento)
persona1.mostrar_datos_personales()
\end{lstlisting}

\subsection{Diagrama de clases}
\clasebox{Persona}
{%
- nombre \\
- apellido \\
- numero\_documento \\
- año\_nacimiento \\
}
{%
+ mostrar\_datos\_personales() \\
}

\newpage
% ---------------------------------------------------------
% EJERCICIO 2.2
% ---------------------------------------------------------
\section{Ejercicio 2.2}

\noindent \textbf{Enlace al repositorio en GitHub:} \\
\url{https://github.com/NatoRemolina/UNAL-Programacion-Orientada-a-Objetos/blob/main/Actividad%202/ejercicio2.py}

\subsection{Código fuente}
\begin{lstlisting}
from enum import Enum

class TipoPlaneta(Enum):
    GASEOSO = "GASEOSO"
    TERRESTRE = "TERRESTRE"
    ENANO = "ENANO"

class Planeta:
    def __init__(self, nombre, cantidad_satelites, masa, volumen, diametro, distancia_sol, tipo, es_observable):
        self.nombre = nombre
        self.cantidad_satelites = cantidad_satelites
        self.masa = masa
        self.volumen = volumen
        self.diametro = diametro
        self.distancia_sol = distancia_sol
        self.tipo = tipo
        self.es_observable = es_observable

    def imprimir(self):
        print(f"Nombre del planeta = {self.nombre}")
        print(f"Cantidad de satélites = {self.cantidad_satelites}")
        print(f"Masa del planeta = {self.masa}")
        print(f"Volumen del planeta = {self.volumen}")
        print(f"Diámetro del planeta = {self.diametro}")
        print(f"Distancia al sol = {self.distancia_sol}")
        print(f"Tipo de planeta = {self.tipo.name}")
        print(f"Es observable = {self.es_observable}")

    def calcular_densidad(self):
        # Evita división por cero
        return self.masa / self.volumen if self.volumen != 0 else 0

    def es_planeta_exterior(self):
        # 149597870 * 3.4
        limite = 508632758.0
        return self.distancia_sol > limite


print("Ingrese el nombre : ")
nombre = input()
print("Ingrese cantidad de satélites : ")
cantidad_satelites = int(input())
print("Ingrese la masa : ")
masa = float(input())
print("Ingrese el volumen : ")
volumen = float(input())

print("Ingrese el diámetro : ")
diametro = int(input())
print("Ingrese la distancia al sol : ")
distancia_sol = float(input())
p1 = Planeta(nombre, cantidad_satelites, masa, volumen, diametro,
             distancia_sol, TipoPlaneta.TERRESTRE, True)
print("********************************")
p1.imprimir()
print(f"Densidad del planeta = {p1.calcular_densidad()}")
print(f"Es planeta exterior = {p1.es_planeta_exterior()}")
print("********************************")
\end{lstlisting}

\subsection{Diagrama de clases}
\clasebox{Planeta}
{%
- nombre \\
- cantidad\_satelites \\
- masa \\
- volumen \\
- diametro \\
- distancia\_sol \\
- tipo \\
- es\_observable \\
}
{%
+ imprimir() \\
+ calcular\_densidad() \\
+ es\_planeta\_exterior() \\
}

\newpage
% ---------------------------------------------------------
% EJERCICIO 2.3
% ---------------------------------------------------------
\section{Ejercicio 2.3}

\noindent \textbf{Enlace al repositorio en GitHub:} \\
\url{https://github.com/NatoRemolina/UNAL-Programacion-Orientada-a-Objetos/blob/main/Actividad%202/ejercicio3.py}

\subsection{Código fuente}
\begin{lstlisting}
class Automovil:
    def __init__(self, marca, modelo, motor, tipo_combustible, tipo_automovil, color, velocidad_maxima, numero_puertas, numero_asientos):
        self.marca = marca
        self.modelo = modelo
        self.motor = motor
        self.velocidad_maxima = velocidad_maxima
        self.numero_puertas = numero_puertas
        self.numero_asientos = numero_asientos

        self.tipo_combustible = tipo_combustible
        self.tipo_automovil = tipo_automovil
        self.color = color

        self._velocidad_actual = 0

    @property
    def tipo_combustible(self):
        return self._tipo_combustible

    @tipo_combustible.setter
    def tipo_combustible(self, valor):
        self._tipo_combustible = valor

    @property
    def tipo_automovil(self):
        return self._tipo_automovil

    @tipo_automovil.setter
    def tipo_automovil(self, valor):
        self._tipo_automovil = valor

    @property
    def color(self):
        return self._color

    @color.setter
    def color(self, valor):
        self._color = valor

    @property
    def velocidad_actual(self):
        return self._velocidad_actual

    def acelerar(self, incremento):
        if incremento < 0:
            print("Error: valor no valido")
            return

        if self._velocidad_actual + incremento <= self.velocidad_maxima:
            self._velocidad_actual += incremento
        else:
            self._velocidad_actual = self.velocidad_maxima
            print(f"Se ha alcanzado la velocidad máxima de {self.velocidad_maxima}")

        print(f"La velocidad actual es {self._velocidad_actual} km/h")

    def desacelerar(self, decremento):
        if decremento < 0:
            print("Error: valor no valido")
            return

        if self._velocidad_actual - decremento >= 0:
            self._velocidad_actual -= decremento
        else:
            self._velocidad_actual = 0
            print("El automóvil se ha detenido por completo.")

        print(f"Velocidad actual: {self._velocidad_actual} km/h")

    def frenar(self):
        self._velocidad_actual = 0
        print("Freno activado. Velocidad: 0 km/h")

    def tiempo_de_llegada(self, distancia):
        if self._velocidad_actual == 0:
            print("No se puede calcular el tiempo de llegada porque el automovil esta detenido.")
            return None

        return distancia / self._velocidad_actual

    def mostrar_atributos(self):
        print("\n--- ESPECIFICACIONES DEL VEHÍCULO ---")
        print(f"Marca: {self.marca} | Modelo: {self.modelo}")
        print(f"Motor: {self.motor}L | Velocidad Máx: {self.velocidad_maxima} km/h")
        print(f"Combustible: {self.tipo_combustible.title()}")
        print(f"Tipo de Auto: {self.tipo_automovil.title()}")
        print(f"Color: {self.color.title()}")
        print(f"Puertas: {self.numero_puertas} | Asientos: {self.numero_asientos}")
        print(f"ESTADO ACTUAL: {self._velocidad_actual} km/h")

def main():
    print("CONFIGURACIÓN DE AUTOMÓVIL\n")

    marca = input("Marca: ")
    modelo = input("Modelo (Año): ")
    motor = float(input("Motor (Litros): "))
    velocidad_maxima = float(input("Velocidad Máxima: "))
    numero_puertas = int(input("Puertas: "))
    numero_asientos = int(input("Asientos: "))

    while True:
        print("\n--- TIPO DE COMBUSTIBLE ---")
        print("1. Gasolina\n2. Bioetanol\n3. Diésel\n4. Biodiésel\n5. Gas Natural")
        opcion_combustible = input("Seleccione (1-5): ")
        if opcion_combustible == "1":
            tipo_combustible_seleccionado = "Gasolina"; break
        elif opcion_combustible == "2":
            tipo_combustible_seleccionado = "Bioetanol"; break
        elif opcion_combustible == "3":
            tipo_combustible_seleccionado = "Diésel"; break
        elif opcion_combustible == "4":
            tipo_combustible_seleccionado = "Biodiésel"; break
        elif opcion_combustible == "5":
            tipo_combustible_seleccionado = "Gas Natural"; break
        else:
            print("Opción inválida.")

    while True:
        print("\n--- TIPO DE AUTO ---")
        print("1. Ciudad\n2. Subcompacto\n3. Compacto\n4. Familiar\n5. Ejecutivo\n6. SUV")
        opcion_tipo = input("Seleccione (1-6): ")
        if opcion_tipo == "1":
            tipo_automovil_seleccionado = "Carro de ciudad"; break
        elif opcion_tipo == "2":
            tipo_automovil_seleccionado = "Subcompacto"; break
        elif opcion_tipo == "3":
            tipo_automovil_seleccionado = "Compacto"; break
        elif opcion_tipo == "4":
            tipo_automovil_seleccionado = "Familiar"; break
        elif opcion_tipo == "5":
            tipo_automovil_seleccionado = "Ejecutivo"; break
        elif opcion_tipo == "6":
            tipo_automovil_seleccionado = "SUV"; break
        else:
            print("Opción inválida.")

    while True:
        print("\n--- COLOR ---")
        print("1. Blanco\n2. Negro\n3. Rojo\n4. Naranja\n5. Amarillo\n6. Verde\n7. Azul\n8. Violeta")
        opcion_color = input("Seleccione (1-8): ")
        if opcion_color == "1":
            color_seleccionado = "Blanco"; break
        elif opcion_color == "2":
            color_seleccionado = "Negro"; break
        elif opcion_color == "3":
            color_seleccionado = "Rojo"; break
        elif opcion_color == "4":
            color_seleccionado = "Naranja"; break
        elif opcion_color == "5":
            color_seleccionado = "Amarillo"; break
        elif opcion_color == "6":
            color_seleccionado = "Verde"; break
        elif opcion_color == "7":
            color_seleccionado = "Azul"; break
        elif opcion_color == "8":
            color_seleccionado = "Violeta"; break
        else:
            print("Opción inválida.")

    mi_auto = Automovil(
        marca,
        modelo,
        motor,
        tipo_combustible_seleccionado,
        tipo_automovil_seleccionado,
        color_seleccionado,
        velocidad_maxima,
        numero_puertas,
        numero_asientos
    )

    while True:
        print("\n--- MENÚ DE CONDUCCIÓN ---")
        print("1. Acelerar")
        print("2. Desacelerar")
        print("3. Frenar")
        print("4. Calcular tiempo de llegada")
        print("5. Ver Ficha Técnica (Atributos)")
        print("6. Salir del programa")

        opcion_menu = input("\nElija una opción: ")

        if opcion_menu == "1":
            cantidad_acelerar = float(input("¿Cuánto desea acelerar? (km/h): "))
            mi_auto.acelerar(cantidad_acelerar)

        elif opcion_menu == "2":
            cantidad_desacelerar = float(input("¿Cuánto desea desacelerar? (km/h): "))
            mi_auto.desacelerar(cantidad_desacelerar)

        elif opcion_menu == "3":
            mi_auto.frenar()

        elif opcion_menu == "4":
            distancia_viaje = float(input("Distancia a recorrer (km): "))
            tiempo_resultado = mi_auto.tiempo_de_llegada(distancia_viaje)
            if tiempo_resultado is not None:
                print(f"Tiempo estimado: {tiempo_resultado:.2f} horas.")

        elif opcion_menu == "5":
            mi_auto.mostrar_atributos()

        elif opcion_menu == "6":
            print("Saliendo del sistema")
            break

        else:
            print("Opción no válida.")


main()
\end{lstlisting}

\subsection{Diagrama de clases}
\clasebox{Automovil}
{%
- marca \\
- modelo \\
- motor \\
- tipo\_combustible \\
- tipo\_automovil \\
- color \\
- velocidad\_maxima \\
- numero\_puertas \\
- numero\_asientos \\
- velocidad\_actual \\
}
{%
+ acelerar() \\
+ desacelerar() \\
+ frenar() \\
+ tiempo\_de\_llegada() \\
+ mostrar\_atributos() \\
}

\newpage
% ---------------------------------------------------------
% EJERCICIO 2.4
% ---------------------------------------------------------
\section{Ejercicio 2.4}

\noindent \textbf{Enlace al repositorio en GitHub:} \\
\url{https://github.com/NatoRemolina/UNAL-Programacion-Orientada-a-Objetos/blob/main/Actividad%202/ejercicio4.py}

\subsection{Código fuente}
\begin{lstlisting}
import math


class Circulo:
    def __init__(self, radio_circulo):
        self.radio = radio_circulo
        self.area = 0
        self.perimetro = 0

    def calcular_area(self):
        self.area = math.pi * math.pow(self.radio, 2)
        return self.area

    def calcular_perimetro(self):
        self.perimetro = 2 * math.pi * self.radio
        return self.perimetro

class Rectangulo:
    def __init__(self, base_rectangulo, altura_rectangulo):
        self.base = base_rectangulo
        self.altura = altura_rectangulo
        self.area = 0
        self.perimetro = 0

    def calcular_area(self):
        self.area = self.base * self.altura
        return self.area

    def calcular_perimetro(self):
        self.perimetro = (2 * self.base) + (2 * self.altura)
        return self.perimetro


class Cuadrado:
    def __init__(self, lado_cuadrado):
        self.lado = lado_cuadrado
        self.area = 0
        self.perimetro = 0

    def calcular_area(self):
        self.area = math.pow(self.lado, 2)
        return self.area

    def calcular_perimetro(self):
        self.perimetro = 4 * self.lado
        return self.perimetro

class Triangulo:
    def __init__(self, base_triangulo, altura_triangulo):
        self.base = base_triangulo
        self.altura = altura_triangulo
        self.area = 0
        self.hipotenusa = 0
        self.perimetro = 0

    def calcular_area(self):
        self.area = (self.base * self.altura)/2
        return self.area

    def calcular_hipotenusa(self):
        self.hipotenusa = math.sqrt(math.pow(self.base, 2) + math.pow(self.altura, 2))
        return self.hipotenusa

    def calcular_perimetro(self):
        self.perimetro = self.base + self.altura + self.calcular_hipotenusa()
        return self.perimetro

    def tipo_triangulo(self):
        # Aseguramos que la hipotenusa esté calculada antes de compararla
        self.calcular_hipotenusa()
        if self.base == self.altura and self.altura == self.hipotenusa:
            return "El triángulo es Equilátero"
        elif self.base == self.altura:
            return "El Triángulo es Isósceles"
        else:
            return "El Triángulo es Escaleno"


def main():
    print("--- datos del circulo ---")
    radio_circulo = float(input("Ingrese el radio del circulo en cm: \n"))
    valor_circulo = Circulo(radio_circulo)

    print("--- Datos Rectangulo ---")
    base_rectangulo = float(input("Ingrese la Base del Rectangulo: \n"))
    altura_rectangulo = float(input("Ingrese la Altura del Rectangulo: \n"))
    valores_rectangulo = Rectangulo(base_rectangulo, altura_rectangulo)

    print("--- Datos Cuadrado ---")
    lado_cuadrado = float(input("Ingrese el Lado del Cuadrado: \n"))
    valores_cuadrado = Cuadrado(lado_cuadrado)

    print("--- datos Triangulo Rectangulo ---")
    base_triangulo = float(input("Ingrese la Base del Triangulo Rectangulo: \n"))
    altura_triangulo = float(input("Ingrese la Altura del Triangulo Rectangulo: \n"))
    valores_triangulo = Triangulo(base_triangulo, altura_triangulo)

    print("--- Resultados de los calculos ---")
    print("--- Círculo ---")
    print(f"Area {valor_circulo.calcular_area():.2f} cm2")
    print(f"Perimetro {valor_circulo.calcular_perimetro():.2f} cm")

    print("--- Rectangulo ---")
    print(f"Area {valores_rectangulo.calcular_area():.2f} cm2")
    print(f"Perimetro {valores_rectangulo.calcular_perimetro():.2f} cm")

    print("--- Cuadrado ---")
    print(f"Area {valores_cuadrado.calcular_area():.2f} cm2")
    print(f"Perimetro {valores_cuadrado.calcular_perimetro():.2f} cm")

    print("--- Triángulo Rectangulo ---")
    print(f"Area {valores_triangulo.calcular_area():.2f} cm2")
    print(f"Perimetro {valores_triangulo.calcular_perimetro():.2f} cm")
    print(f"Hipotenusa {valores_triangulo.calcular_hipotenusa():.2f} cm")
    print(f"Tipo {valores_triangulo.tipo_triangulo()}")

main()
\end{lstlisting}

\subsection{Diagrama de clases}
\clasebox{Circulo}
{%
- radio \\
}
{%
+ calcular\_area() \\
+ calcular\_perimetro() \\
}

\clasebox{Rectangulo}
{%
- base \\
- altura \\
}
{%
+ calcular\_area() \\
+ calcular\_perimetro() \\
}

\clasebox{Cuadrado}
{%
- lado \\
}
{%
+ calcular\_area() \\
+ calcular\_perimetro() \\
}

\clasebox{Triangulo}
{%
- base \\
- altura \\
- hipotenusa \\
}
{%
+ calcular\_area() \\
+ calcular\_perimetro() \\
+ calcular\_hipotenusa() \\
+ tipo\_triangulo() \\
}

\newpage
% ---------------------------------------------------------
% EJERCICIO 2.5
% ---------------------------------------------------------
\section{Ejercicio 2.5}

\noindent \textbf{Enlace al repositorio en GitHub:} \\
\url{https://github.com/NatoRemolina/UNAL-Programacion-Orientada-a-Objetos/blob/main/Actividad%202/ejercicio5.py}

\subsection{Código fuente}
\begin{lstlisting}
class CuentaBancaria:
    def __init__(self, nombre, apellido, id, tipo_de_cuenta):
        self.nombre = nombre
        self.apellido = apellido
        self.id = id
        self.tipo_de_cuenta = tipo_de_cuenta
        self.saldo = 0

    def consignar(self, valor_consignado):
        self.saldo = self.saldo + valor_consignado
        print(f"Consignación exitosa. Nuevo saldo: {self.saldo}")

    def retirar(self, valor_retirado):
        if self.saldo >= valor_retirado:
            self.saldo = self.saldo - valor_retirado
            print(f"Retiro exitoso. Saldo actual: {self.saldo}")
        else:
            print("Saldo insuficiente")

    def imprimir_datos(self):
        print(f"Nombre del titular: {self.nombre}")
        print(f"Apellido del titular: {self.apellido}")
        print(f"Identificación: {self.id}")
        print(f"Tipo de cuenta: {self.tipo_de_cuenta}")
        print(f"Saldo de la cuenta: {self.saldo}")

def main():
    print("--- Ingrese los datos ---")
    nombre = input("Nombres del titular: \n")
    apellido = input("Apellidos del titular: \n")
    id = input("Número de cuenta: \n")

    while True:
        tipo_de_cuenta = input("Tipo de cuenta: \n").lower()
        if tipo_de_cuenta == "ahorro" or tipo_de_cuenta == "corriente":
            break
        else:
            print("Tipo de cuenta no valido.")

    cuenta = CuentaBancaria(nombre, apellido, id, tipo_de_cuenta)

    while True:
        print("--- menu ---")

        print("--- 1. consignar valor ---")
        print("--- 2. retirar valor ---")
        print("--- 3. consultar datos y saldo ---")
        print("--- 4. salir ---")

        opcion = input("Ingrese la opcion: \n")

        if opcion == "1":
            valor_consignar = float(input("Ingrese el valor a consignar: \n"))
            cuenta.consignar(valor_consignar)
        elif opcion == "2":
            valor_retirar = float(input("Ingrese valor a retirar: \n"))
            cuenta.retirar(valor_retirar)
        elif opcion == "3":
            cuenta.imprimir_datos()
        elif opcion == "4":
            print("Finalizado exitosamente")
            break
        else:
            print("Opcion no valida")

main()
\end{lstlisting}

\subsection{Diagrama de clases}
\clasebox{CuentaBancaria}
{%
- nombre \\
- apellido \\
- id \\
- tipo\_de\_cuenta \\
- saldo \\
}
{%
+ consignar() \\
+ retirar() \\
+ imprimir\_datos() \\
}

\end{document}
